% English translation, made from our own transcription (original.tex), not from any existing
% translation. Every math expression, symbol, \tag, \origpage{N} marker and \emph run is reproduced
% unchanged; only the prose is translated.
%
% PROSE INSIDE FORMULA TRANSLATED (HOUSESTYLE R21):
% In the formula on p. 1146 ($u = \text{const.}$), \text{const.} is the standard abbreviation in both
% French and English and is preserved unchanged.
%
% QUOTATION MARKS ARE NOT TRANSLATED (HOUSESTYLE R18, R22):
% The paragraph-initial « / » of the Comptes rendus and the terminal » at the end of p. 1147 are
% typographic features of the source edition and stand exactly where the transcription puts them,
% set tight against the text.
%
% FAITHFULNESS TO APPARENT PRINTER'S ERRORS (HOUSESTYLE R4):
% - p. 1146, display 1: x^{2}(az^{2} + 2b - z + c). The print sets a dash/hyphen between b and z.
%   Reproduced faithfully and noted with \ednote.
% - p. 1146, line 11: "deux polygones de degré 2 et 4 en z" is translated faithfully as
%   "two polygons of degree 2 and 4 in z", preserving the printer's slip ("polygons" for
%   "polynomials") and noted with \ednote.
% - p. 1146, line 14: "» Il est ailleurs aisé de voir" (print sets "ailleurs" for "d'ailleurs")
%   translated naturally into English prose ("It is, moreover, easy to see").
% - p. 1146, line 18: "» De inême" (print sets "inême" for "même") translated naturally into English
%   prose ("Likewise").
%
% TERMINOLOGY:
%   * intégrales de différentielles totales = "integrals of total differentials"
%   * première espèce = "first kind"
%   * surfaces du quatrième ordre = "surfaces of the fourth order"
%   * surface réglée = "ruled surface"
%   * surface de révolution = "surface of revolution"
%   * unicursale = "unicursal"
%   * point conique du second ordre = "conical point of the second order"
%   * cône tangent indécomposable = "indecomposable tangent cone"
%   * polynômes entiers = "polynomials"
%   * théorème d'Abel = "Abel's theorem"
\documentclass{article}
\usepackage{readmasters}
\begin{document}

\origpage{1145}
\section*{Mathematical Analysis. --- On integrals of total differentials. Note by M.~H.~Poincaré, presented by M.~Hermite.}

«The recent discovery of M.~Picard on total differentials of the first kind has opened to analysts an entirely new path, where they will undoubtedly encounter many important propositions. It will therefore perhaps not be without interest to point out here certain partial results which, although very easy to prove, may be useful to geometers who concern themselves with this question.

\origpage{1146}
»I sought first to determine which surfaces of the fourth order possess integrals of the first kind. I found that all these surfaces can be reduced, by a linear change of variables, either to the ruled surface
\[
\begin{aligned}
x^{2}(az^{2} + 2b - z + c) &+ 2xy(a'z^{2} + 2b'z + c') \\
&+ y^{2}(a''z^{2} + 2b''z + c'') = 0,
\end{aligned}
\]\ednote{Printed $2b - z + c$ --- an apparent misprint for $2bz$, matching $2b'z$ and $2b''z$, and the denominator of the integral immediately following. Reproduced here as printed.}
which admits the integral
\[
\int \frac{y\,dz}{x(az^{2} + 2bz + c) + y(a'z^{2} + 2b'z + c')},
\]
or to the surface of revolution
\[
(x^{2} + y^{2})^{2} + 2(x^{2} + y^{2})Z_{2} + Z_{4} = 0
\]
($Z_{2}$ and $Z_{4}$ denoting two polygons\ednote{Printed ``polygones'' --- an apparent compositor's slip for ``polynômes''. Reproduced here as printed.} of degree 2 and 4 in $z$), which admits the integral
\[
\int \frac{dz}{x^{2} + y^{2} + Z_{2}}.
\]

»It is, moreover, easy to see that all ruled surfaces and all surfaces of revolution admit integrals of the first kind, unless, of course, they are unicursal.

»If a surface admits an integral of the first kind $u$, reducible to elliptic integrals, the curves $u = \text{const.}$ are algebraic.

»If one can trace on a surface a unicursal curve, and if $u$ is an arbitrary integral of the first kind of this surface, the value of $u$ will be the same all along the curve.

»Likewise, if the surface admits a conical point of the second order, whose tangent cone is indecomposable, the value of $u$ at this conical point will be \emph{determinate}.

»If one can trace on a surface two series of unicursal curves, it will have no integral of the first kind; if, on a non-unicursal surface, one can trace a series of unicursal curves, in such a way that through each point of the surface there passes, in general, only a single one of these curves, it will have integrals of the first kind.

»Suppose that a surface is generated by the elimination of two
\origpage{1147}
parameters $a$ and $b$, between the three equations
\[
\begin{gathered}
\varphi(x, y, z, a, b) = 0, \\
\varphi_{1}(x, y, z, a, b) = 0, \\
\psi(a, b) = 0.
\end{gathered}
\]

»If the three polynomials $\varphi$, $\varphi_{1}$ and $\psi$ are \emph{the most general of their degrees}, the relation $\psi = 0$ is of genus greater than 0; to a point of the surface corresponds a single system of values of the parameters and, consequently, the surface admits integrals of the first kind.

»Finally, Abel's theorem applies to integrals of total differentials.

»Let $M_{1}, M_{2}, \ldots, M_{q}$ be the points of intersection of the surface with the curve
\[
\frac{\alpha}{\lambda} = \frac{\beta}{\mu} = \frac{\gamma}{\nu},
\]
$\alpha$, $\beta$, $\gamma$ being polynomials in $x$, $y$, $z$ and $\lambda$, $\mu$, $\nu$ constants. Let $u_{1}, u_{2}, \ldots, u_{q}$ be the values of a certain integral of the first kind $u$ at these different points.

»Let $M'_{1}, M'_{2}, \ldots, M'_{q}$ be the points of intersection of the surface with the curve
\[
\frac{\alpha}{\lambda'} = \frac{\beta}{\mu'} = \frac{\gamma}{\nu'},
\]
$\lambda'$, $\mu'$ and $\nu'$ being new constants. Let $u'_{1}, u'_{2}, \ldots, u'_{q}$ be the values of the integral $u$ at the points $M'_{1}, M'_{2}, \ldots, M'_{q}$.

»We will have
\[
u_{1} + u_{2} + \ldots + u_{q} = u'_{1} + u'_{2} + \ldots + u'_{q},
\]»

\end{document}
